% =============================================================================
% Section 6: Implications
% =============================================================================
\section{Implications}\label{sec:implications}

The duplex finding has consequences for both policy and measurement.
This section states them, bounded by the evidence.

\subsection{Policy implications}

\subsubsection{Defence supply-chain diversification is the
  highest-leverage intervention}

If a government seeks to reduce its {\ISI} composite score---that is,
to reduce the concentration of its bilateral external
dependencies---the highest-return intervention is to diversify
defence procurement.  A one-standard-deviation reduction in a
country's defence axis score produces a larger effect on the composite
than an equivalent reduction in any other axis, precisely because
the defence axis contributes the most composite variance and has
the widest cross-country dispersion.

This is not a normative recommendation.  Defence procurement is
constrained by alliance structures, technology access, and
interoperability requirements that this index does not measure.  The
point is mechanical: the index rewards defence diversification more
than it rewards diversification in any other domain.

\subsubsection{Logistics diversification has similar ranking power}

Logistics contributes nearly as much variance as defence.  Countries
dominated by the logistics axis (11 of~27) can reduce their composite
score most efficiently by diversifying port and freight
dependencies.  For small island economies and peripheral states
where logistics concentration is structurally determined by
geography, the implication is that the index will persistently
rank them as more concentrated regardless of reforms in other domains.

\subsubsection{Financial and energy reforms are ranking-neutral}

Reforms that reduce financial or energy axis scores will reduce the
country's composite score arithmetically, but will not change the
country's rank position.  Financial contributes 1.0\% of composite
variance; energy contributes 3.5\%.  The ranking is insensitive to
variation in these domains.

This does not mean financial or energy dependence is unimportant.
It means the {\ISI} composite ranking does not differentiate
countries along these dimensions.

\subsection{Measurement implications}

\subsubsection{Data quality in defence and logistics is critical}

The two axes that drive the ranking rely on data sources that are,
respectively, sparse and opaque:

\begin{itemize}[leftmargin=2em]
  \item \textbf{Defence.}  The {\ISI} defence axis is constructed
    from SIPRI trend-indicator values~\cite{sipri_transfers}, which
    cover major conventional arms transfers.  Small arms, dual-use
    items, cyber capabilities, and maintenance contracts are excluded.
    Countries with highly classified procurement programmes may have
    understated concentration scores.

  \item \textbf{Logistics.}  The logistics axis aggregates
    containerised freight flows and port throughput data.  The
    underlying data are collected by national statistical offices and
    Eurostat~\cite{eurostat_comext} with varying granularity.  Transit
    flows, overland freight, and air cargo are not included in the
    current construction.
\end{itemize}

\noindent
Any measurement error in these two domains propagates directly to
the composite ranking.  Errors in financial or energy data are,
by contrast, ranking-irrelevant.

\subsubsection{The equal-weight assumption amplifies
  high-variance axes}

The equal-weight mean assigns each axis the same nominal weight.  But
effective weight---the share of composite variance attributable to
each axis---is determined by the empirical dispersion.  Under equal
weights, defence and logistics receive $1/6 \approx 16.7\%$
nominal weight each but contribute $35.1\%$ and $34.3\%$ of the
variance, respectively.

This is not a flaw.  It is an inherent property of linear
aggregation with heterogeneous sub-indicators.  Users who wish to
equalise effective variance shares would need to rescale axis scores
(e.g., z-score standardisation) before averaging.
Paper~2~\cite{isi_paper2} reports that z-score standardisation
produces near-identical rankings, suggesting the result is not an
artefact of the scale.

\subsection{Boundaries}

The implications stated above are bounded by the following constraints:

\begin{itemize}[leftmargin=2em]
  \item They apply to the EU-27 under vintage-2024 data.  The
    variance structure may differ for other country sets or time
    periods.
  \item They concern the \emph{ranking} produced by the index, not
    the underlying economic or strategic significance of each
    domain.  A country may face severe financial concentration risk
    that the ranking does not differentiate.
  \item They do not imply that defence or logistics should receive
    higher (or lower) nominal weight.  The observation is
    descriptive, not prescriptive.
\end{itemize}
